% =====================================================================
%  Part VII -- Numerical analysis prerequisites.
%  Source: tier0.md Part V (2483-2848), order unchanged.
%
%  NOTE: tier0.md:2569-2570 is a FALSE markdown blockquote -- a
%  line-wrapped ">" sign in "cond(S_active) < 1e6 pass / > 1e8 fail".
%  It is set inline here, not as a callout.
% =====================================================================
\part{Numerical analysis prerequisites}
\label{part:numerics}

% ---------------------------------------------------------------------
\chapter{Catastrophic cancellation --- \texorpdfstring{$\kappa$}{kappa} is a
  condition number}
\label{sec:cancellation}

The single most important piece of numerical analysis in the project, and it
deserves a derivation rather than an assertion.

\section{The derivation}
\label{sec:cancellationderivation}

Compute $d = a - b$ where $a \approx b > 0$, with each input carrying relative
error $\varepsilon$ (the standard rounding model ---
\citealt{goldberg1991}; \citealt{higham2002}):

\begin{equation*}
\mathrm{fl}(a) = a(1+\varepsilon_1), \qquad
\mathrm{fl}(b) = b(1+\varepsilon_2), \qquad
|\varepsilon_i| \leq \varepsilon
\end{equation*}

\begin{equation*}
\text{computed } d = (a-b) + (a\varepsilon_1 - b\varepsilon_2)
\end{equation*}

Relative error of the result:

\begin{equation}
\frac{|\Delta d|}{|d|} \;\leq\; \varepsilon\,\frac{a+b}{|a-b|}
\label{eq:cancellation-bound}
\end{equation}

Now recognize the denominator. With $f^+ = a$ and $f^- = b$:

\begin{equation}
\kappa \equiv \frac{|f^+ - f^-|}{f^+ + f^-} \qquad \implies \qquad
\frac{|\Delta d|}{|d|} \leq \frac{\varepsilon}{\kappa}
\label{eq:kappa-condition}
\end{equation}

\begin{result}
\textbf{$\kappa$ is exactly the reciprocal condition number of the
subtraction.} (The condition number of $a - b$ is $(|a|+|b|)/|a-b|$ ---
\citealt{higham2002}; for positive $f^\pm$ that is exactly $1/\kappa$.) It is
not a heuristic equilibrium diagnostic that happens to be useful. It is
\emph{the} error-amplification factor for recovering net flux from gross
quantities.
\end{result}

Note that \S\ref{sec:detailedbalance} arrived at $\kappa$ from a completely
different direction --- as the numerical test of detailed balance,
eq.~\eqref{eq:kappa-detailed-balance}. The two meanings coincide, which is why
$\kappa$ is load-bearing in both the physics and the numerics of this project.

\section{The numbers that follow}
\label{sec:kappanumbers}

\begin{center}
\small
\begin{tabularx}{\textwidth}{@{}l l R{1.125} R{0.875}@{}}
\toprule
\textbf{$\kappa$} & \textbf{Amplification $1/\kappa$} &
\textbf{Net error from float32 gross ($\varepsilon = 1.2\times10^{-7}$)} &
\textbf{float64 ($\varepsilon = 2.2\times10^{-16}$)} \\
\midrule
1          & 1          & $1.2\times10^{-7}$ & $2.2\times10^{-16}$ \\
$10^{-1}$  & 10         & $1.2\times10^{-6}$ & $2.2\times10^{-15}$ \\
$10^{-3}$  & $10^{3}$   & $1.2\times10^{-4}$ & $2.2\times10^{-13}$ \\
$10^{-6}$  & $10^{6}$   & 0.12 --- \textbf{destroyed} & $2.2\times10^{-10}$ \\
$10^{-9}$  & $10^{9}$   & --- & $2.2\times10^{-7}$ \\
\bottomrule
\end{tabularx}
\end{center}

\textbf{Now the project's thresholds are derivable rather than arbitrary:}

\begin{itemize}
\item The \textbf{$\kappa > 0.1$ active-set gate} is the statement
      ``amplification $\leq 10\times$'' --- one decade of precision loss,
      tolerable.
\item \textbf{float64 everywhere in conservation arithmetic} buys the headroom:
      the same $\kappa = 10^{-6}$ that would be fatal in float32 is harmless in
      float64.
\item The \textbf{$1/\kappa$ error amplification if $\varphi$ were inferred from
      gross-scale features} --- the phrasing used in the verdict document --- is
      precisely \eqref{eq:kappa-condition}.
\item And the reason the worst-species coverage failure in
      $\Tn \in [4.0, 6.3)$ matters: those species' net evolution rides on
      columns with $\kappa \in 10^{-3}$--$10^{-1}$, i.e.\ 10--1000$\times$
      amplification.
\end{itemize}

\section{Why this reframes the whole kill-test}
\label{sec:killtestreframe}

The kill-test asks: \emph{can a model that predicts gross-scale quantities
recover the net?} Equation \eqref{eq:kappa-condition} says the answer is
governed entirely by the $\kappa$ distribution. The kill-test is therefore a
\textbf{numerical conditioning study wearing a physics costume} --- and that is
why its thresholds are stated in $\kappa$ and in $\cond(S_{\mathrm{active}})$,
both condition numbers.

% ---------------------------------------------------------------------
\chapter{Conditioning and the rank-revealing subtlety}
\label{sec:conditioning}

For a linear system $A x = b$, relative error amplification is the condition
number

\begin{equation*}
\cond(A) = \frac{\sigma_{\max}}{\sigma_{\min}} \qquad
(\text{ratio of extreme singular values})
\end{equation*}

\textbf{The project's use.} Given $\dd Y$, recover $\varphi$ from
$\dd Y = \nu\varphi$. With $n_{\text{species}} \approx 80$--151 equations and
$n_{\text{reactions}} \approx 607$--1518 unknowns this is massively
\textbf{underdetermined} --- infinitely many $\varphi$ produce the same $\dd Y$,
and the minimum-norm solution is $\varphi = \nu^{+}\,\dd Y$ via the
Moore--Penrose pseudo-inverse. The sensitivity of that recovery to perturbations
in $\dd Y$ is set by the condition number of the relevant submatrix. Hence the
$\cond(S_{\mathrm{active}}) < 10^{6}$ pass / $> 10^{8}$ fail gate, where
$S_{\mathrm{active}}$ is $\nu$ restricted to the net columns that are actually
carrying flux.

Underdetermination is not itself a problem --- it is \emph{why} Target A works,
since any $\varphi$ in the null space decodes to the same conserving $\dd Y$.
The problem would be ill-conditioning within the active set, which is what the
gate measures.

\textbf{The subtlety you must not get wrong.} $\nu$ has an \emph{exact
structural left-null vector}: $A\cdot\nu = 0$, by baryon conservation, to
machine precision. So $\sigma_{\min} = 0$ identically, and a naive
\code{np.linalg.cond} returns ${\sim}10^{16}$ --- which is not
ill-conditioning, it is the conservation law showing up as a zero singular
value.

The meaningful quantity is the \textbf{rank-revealing} condition number --- the
ratio of extreme \emph{nonzero} singular values, i.e.\ restricted to the row
space (the name borrows from rank-revealing factorizations,
\citealt{chan1987}; ``effective condition number'' is the closest standard
term). That gives the measured 41.57 / 57.86 --- perfectly well-conditioned.

\begin{warn}
\textbf{Say which matrix, because the nullity differs.} On the nuclei-only $\nu$
($n \times r$) the only exact left-null vector is $A$, so there is \textbf{one}
zero singular value to drop --- $Z$ is \emph{not} one, because
$Z\cdot\nu_j = d_{\mathrm{electron},j} \neq 0$ on weak columns
(\S\ref{sec:ye}). On the extended
$\tilde\nu = \mathrm{vstack}(\nu, \text{lepton ledgers})$
($(n+3) \times r$), which is what \code{graph/stoich.py} builds and validates,
\textbf{every row of $C$ annihilates it}: $C\tilde\nu = 0$ exactly for all three
independent rows, so $\tilde\nu$ has a \textbf{three}-dimensional exact
left-null space and a rank-revealing $\cond$ must discard three singular values,
not one. Dropping the wrong number silently reports either a spurious $10^{16}$
or a spurious well-conditioning. The 41.57 / 57.86 figures are the nuclei-only
reading (\code{RESULTS.md} 2026-07-09, post-reconciliation, superseding the
earlier 41.70 / 57.55).
\end{warn}

\begin{aside}
A prior session in this repo rendered $\cond(\nu) = 1.5\times10^{16}$ in a
figure whose caption claimed $\approx 42$. It was caught in review. The lesson
generalizes: \textbf{an exact structural null space is not ill-conditioning},
and any condition number computed on a matrix with known exact null vectors must
be rank-revealing.
\end{aside}

% ---------------------------------------------------------------------
\chapter{Stiffness, defined properly}
\label{sec:stiffness}

``Stiff'' is often used loosely to mean ``fast''. It means something specific.

Let $J = \nu\,\partial R/\partial Y$ be the Jacobian of the network ODE, with
eigenvalues $\lambda_i$ (mostly with $\mathrm{Re}\,\lambda_i < 0$ --- decaying
modes). Define the \textbf{stiffness ratio} (\citealt{hix1999b}, after Lambert)

\begin{equation}
S = \frac{\max_i |\mathrm{Re}\,\lambda_i|}{\min_i |\mathrm{Re}\,\lambda_i|}
\label{eq:stiffness-ratio}
\end{equation}

An explicit method is stable only for step $h \lesssim 2/\max|\lambda|$ ---
forward-Euler stability, observed to the fourth digit in a real network
(``precisely two over the fastest rate'' --- \citealt{guidry2013}). Integrating
out to the \emph{slowest} timescale $t \sim 1/\min|\lambda|$ therefore requires

\begin{equation*}
N_{\text{steps}} \sim S
\end{equation*}

At silicon-burning conditions: fastest modes are photodisintegration /
$(n,\gamma)$-$(\gamma,n)$ pairs with rates up to ${\sim}10^{10}$~s$^{-1}$;
slowest are weak reactions at ${\sim}10^{-5}$--$10^{-2}$~s$^{-1}$ (both
endpoints are order-of-magnitude estimates --- \emph{assumed}, pending the
Jacobian measurement of self-check item~6 --- but the bracket they produce sits
inside the literature envelope: ``$S > 10^{15}$ is not uncommon in
astrophysics'', \citealt{hix1999b}; fastest-to-slowest spans of 10--20 orders,
\citealt{guidry2013}). So

\begin{equation*}
S \sim 10^{12}\text{--}10^{15}
\end{equation*}

Explicit integration would need ${\sim}10^{15}$ steps to follow the \Ye{}
evolution. \textbf{Implicit is not an optimization, it is a necessity.}

\section{Stiffness and QSE are the same phenomenon}
\label{sec:stiffnessqse}

The fast modes decay quickly, after which the trajectory lies on a
low-dimensional \textbf{slow manifold} in composition space. On that manifold
the fast reactions are in balance --- which is exactly the definition of
quasi-statistical equilibrium.

\begin{result}
\textbf{Stiffness $\iff$ the existence of a slow manifold $\iff$ QSE.} S9
(equilibrium theory) and S10 (stiff integration) are two views of one structure.
Reading them as unrelated topics is the main way people fail to understand this
domain. (Strictly the arrows run one way --- stiffness alone guarantees a fast
decaying mode, not a physically meaningful equilibrated sector; the
identification is the standard singular-perturbation picture, cf.\ ILDM ---
\citealt{maas1992}.)
\end{result}

Note this also connects back to \S\ref{sec:onezone}: the slow manifold \emph{is}
the reachable manifold, approached after transients decay. The three concepts
--- stiffness, QSE, and the reachable set --- are one object seen from numerics,
thermodynamics, and sampling respectively.

This also explains why the \emph{genuine} unlock for hot-state integration cost
--- as the project's own negative-results investigation concluded --- is
\textbf{QSE-reduced integration}: algebraically eliminate the fast equilibrated
sector and integrate only the slow one. Not a solver trick; exploiting the slow
manifold directly.

% ---------------------------------------------------------------------
\chapter{Floating point, concretely}
\label{sec:floatingpoint}

\begin{center}
\begin{tabular}{@{}l l l@{}}
\toprule
\textbf{Type} & \textbf{$\varepsilon$ (machine epsilon)} &
\textbf{Decimal digits} \\
\midrule
float32 & $1.19\times10^{-7}$  & ${\sim}7$ \\
float64 & $2.22\times10^{-16}$ & ${\sim}16$ \\
\bottomrule
\end{tabular}
\end{center}

\textbf{Why the conservation gate is achievable.} Naive summation of $n$ terms
accumulates error ${\sim}n\varepsilon$ in the worst case (the sharp bound is
$(n-1)u\cdot\sum|x_i| + O(u^2)$ --- \citealt{higham2002}; $n\varepsilon$ at
$O(1)$ gross scale is the fair simplification). For the largest network:

\begin{equation*}
n\,\varepsilon = 151 \times 2.22\times10^{-16} \approx 3.4\times10^{-14}
\end{equation*}

comfortably below the $10^{-12}$ gate --- about 1.5 decades of margin. In
float32 the same sum gives $1.8\times10^{-5}$, missing the gate by seven orders
of magnitude. \textbf{This is the derivation behind ``float64 for anything
touching conservation.''}

\textbf{Why the drift bound is scaled rather than absolute.} An absolute bound of
$10^{-12}$ is meaningless when $|\varphi| \sim 10^{-20}$ --- every quantity
involved is already below it, so the test would pass vacuously. Hence the
project's bound

\begin{equation*}
\max\!\bigl(10^{-12}\,s,\; 10^{-13}\,G\bigr), \qquad
s = \min(1, \max|\varphi|), \qquad
G = |A|\cdot(|\nu|\cdot|\varphi|)
\end{equation*}

which reduces to the hard absolute gate at $O(1)$ magnitudes and stays
machine-precision-tight relative to the \emph{gross} scale at the small end.
Read \code{tests/test\_conservation.py} with this in mind --- the tolerance
design is the interesting part of that file.

(In practice the measured column drifts are exactly 0.0, because $\nu$ is
integer and the products are exactly representable. The bound exists for the
general random-$\varphi$ case.)

% ---------------------------------------------------------------------
\chapter{Newton's method, damping, and globalization}
\label{sec:newton}

Prerequisite for reading \code{qse/solver.py}, whose structure looks defensive
until you know what it is defending against.

\section{The method}
\label{sec:newtonmethod}

To solve $\mathbf{F}(\mathbf{u}) = \mathbf{0}$, iterate

\begin{equation}
\mathbf{u}_{k+1} = \mathbf{u}_k - J^{-1}F(\mathbf{u}_k), \qquad
J = \frac{\partial F}{\partial \mathbf{u}}
\label{eq:newton}
\end{equation}

Near a simple root, convergence is \textbf{quadratic}: the error squares each
step, so $10^{-2} \to 10^{-4} \to 10^{-8} \to 10^{-16}$. That is why Newton is
the default whenever an analytic Jacobian is available --- and here it is,
because the Saha form differentiates in closed form.

\section{Why it fails, specifically here}
\label{sec:newtonfails}

The NSE residual involves $X_i = \exp(\log C_i + (Z_i u_p + N_i u_n)/kT)$. With
$Z, N \sim 28$ and $kT \sim 0.3$~MeV, the exponent has a coefficient of order
100 per MeV of chemical potential. A Newton step that overshoots by even a few
MeV pushes the exponent by hundreds --- producing \code{inf}, then \code{nan},
then a dead iteration.

Two guards, both visible in the code:

\begin{enumerate}
\item \textbf{Step damping.} Cap the step, $\|\delta u\| \leq 2$~MeV per
      iteration. This sacrifices quadratic convergence far from the root in
      exchange for never leaving the representable region. Standard
      globalization; a line search would be the fancier alternative.
\item \textbf{Exponent clipping} at \code{EXP\_CLIP} (mirroring pynucastro), so
      an overshoot saturates rather than overflowing.
\end{enumerate}

\section{Globalization: the fallback that cannot fail}
\label{sec:globalization}

Damped Newton is still only \emph{locally} convergent. \code{solve\_nse}
therefore falls back to \textbf{nested bisection}: slow (linear convergence, one
bit per iteration) but \emph{guaranteed} on a bracketed monotone problem. The
residual is monotone in each chemical potential, which is what makes bracketing
valid.

Convergence tolerance is \code{tol = 1e-11} on the residual.

\begin{aside}
\textbf{The pattern is worth extracting, because the repo uses it repeatedly:} a
fast method, a provably convergent fallback, and a \textbf{record of which one
ran} --- \code{NSEResult.method} is \code{"newton"} or \code{"bisection"}.
Silent fallback would hide a systematic problem; recorded fallback turns it into
a measurable rate. The batched solvers added later preserve exactly this:
vectorized Newton with per-row convergence masking, delegating non-converged
rows to the proven scalar path.
\end{aside}

% ---------------------------------------------------------------------
\chapter{ODE error control: what rtol and atol actually promise}
\label{sec:errorcontrol}

Prerequisite for reading \code{fluxes/integrate.py}.

\section{Local versus global error}
\label{sec:localglobal}

An adaptive solver controls the \textbf{local} error per step against

\begin{equation}
\text{tolerance}_i = \mathrm{rtol}\cdot|y_i| + \mathrm{atol}_i
\label{eq:error-tolerance}
\end{equation}

(\citealt{hairer1993}; SciPy's \code{solve\_ivp} documents exactly this contract
--- ``keeps the local error estimates less than
$\mathrm{atol} + \mathrm{rtol}\cdot\mathrm{abs}(y)$''.) It does \textbf{not}
control global error. The relationship between them is exactly
\S\ref{sec:rollouterror}'s accumulation problem applied to the integrator:
global error is the accumulation of local errors under the flow's own
contraction. For a contracting (dissipative) problem --- which this is --- local
control is a good proxy. For a chaotic or neutrally stable one it would not be.

\textbf{So \code{rtol = 1e-8} does not mean ``the answer is right to 1e-8''.} It
means each step's local error estimate was held there. Verifying the actual
accuracy needs either a tolerance sweep or an independent method.

Both are done here: the integrator is checked BDF versus Radau (a multistep
method versus an implicit Runge--Kutta --- different error structures entirely)
and across \code{rtol} from \code{1e-10} to \code{5e-8}. Agreement across
\emph{method} and \emph{tolerance} is the real accuracy evidence; the tolerance
number alone is not.

\section{Why atol matters more than usual}
\label{sec:atol}

With a state spanning 15 decades, $\mathrm{rtol}\cdot|y|$ is meaningless for
trace species: a species at $Y \sim 10^{-20}$ would have a tolerance of
$10^{-28}$, and the solver would take absurdly small steps chasing numerical
noise. \textbf{atol sets the ``I do not care below this'' floor}, and choosing it
is a physics decision, not a numerical one.

The integrator uses \code{atol\_y = atol\_phi = 1e-18}, and --- the point worth
noting --- carries them as \textbf{separate parameters}. The augmented state is
$[\mathbf{Y}, \boldsymbol{\Phi}]$, and abundances and cumulative fluxes have
different natural scales; a single atol would impose one species' notion of
negligible on the other.

\section{Dense output and t\_eval}
\label{sec:denseoutput}

\code{t\_eval} requests values at specified times. The solver does \textbf{not}
change its step sequence to land on them --- it interpolates using the
polynomial it already built for error estimation. This is why one integration to
$\dd t = 10^{2}$~s with \code{t\_eval} at all nine label $\dd t$s costs
essentially the same as integrating to $10^{2}$~s alone, and gives all nine
labels. That is the design that makes the label prototype affordable at all.

\section{BDF versus Radau}
\label{sec:bdfradau}

\begin{itemize}
\item \textbf{BDF} --- multistep, variable order 1--5, cheap per step (one
      Jacobian reused across steps), the standard choice for large stiff
      systems.
\item \textbf{Radau} --- implicit Runge--Kutta (Radau IIA, order 5), L-stable,
      more robust on very stiff or strongly nonlinear problems, more expensive
      per step.
\end{itemize}

(\citealt{hairer1996}; properties confirmed against the SciPy implementation.
BDF above order 2 is only A($\alpha$)-stable --- the Dahlquist barrier --- which
is precisely why an L-stable Radau cross-check is meaningful.)

Using BDF as production and Radau as cross-check is the right split: you get
throughput from one and solver-independence evidence from the other.

% ---------------------------------------------------------------------
\chapter*{Self-check --- numerical analysis}
\addcontentsline{toc}{chapter}{Self-check --- numerical analysis}
\label{sec:selfcheck-numerics}

\begin{selfcheck}
\item Derive \eqref{eq:cancellation-bound} carefully, tracking signs.
\item At what $\kappa$ does float64 gross-flux arithmetic fail to deliver the
      $3\times10^{-6}$ per-step \Ye{} gate? Compare with the measured $\kappa$
      distribution.
\item Target A predicts $\varphi$ directly rather than differencing $f^+$ and
      $f^-$. Does that escape \eqref{eq:kappa-condition}? Argue both ways ---
      then note what the labels themselves are made of.
\item Explain why $\dd Y = \nu\varphi$ is underdetermined, why that is \emph{not}
      a problem for Target A, and what would actually constitute a conditioning
      failure.
\item Compute \code{np.linalg.cond} on an exported $\nu$ and reproduce the
      ${\sim}10^{16}$ artifact. Then compute the rank-revealing version and
      recover $\approx 42$ / $\approx 58$.
\item Estimate $S$ from a real Jacobian at $\Tn = 4$ and compare with the
      $10^{12}$--$10^{15}$ claim. Which reactions set each end?
      \label{sc:jacobian}
\item Show that the scaled drift bound reduces to the absolute $10^{-12}$ gate at
      $|\varphi| \sim 1$, and explain why an absolute bound alone would be
      vacuous at $10^{-20}$.
\item Estimate how far a single undamped Newton step can move the Saha exponent
      if $\delta u = 10$~MeV, at $kT = 0.345$~MeV and $Z = N = 28$. Now justify
      the 2~MeV cap.
\item Why is bisection a \emph{valid} fallback here --- what property of the
      residual does bracketing require, and does the QSE residual have it too?
\item State precisely what \code{rtol = 1e-8} promises and what it does not. Then
      describe the two independent checks that would establish actual accuracy.
\item Why does the augmented integrator carry separate \code{atol\_y} and
      \code{atol\_phi}? Construct a state where a single shared atol would be
      wrong for one of them.
\end{selfcheck}

\begin{repopointer}
\textbf{Code anchors:} \code{tests/test\_conservation.py} (tolerance design),
\code{graph/metrics.py::condition\_numbers} and \code{drift\_metrics},
\code{killtest/active\_set.py::cond\_s\_active}, \code{fluxes/engine.py} (where
$\kappa$ is formed), \code{qse/solver.py} (damped Newton $+$ bisection,
\code{NSEResult.method}), \code{fluxes/integrate.py} (BDF/Radau, rtol/atol,
\code{t\_eval}).
\end{repopointer}
